\subsection{Aging Assessment: Vibration-Induced Fatigue}
\label{sec:aging}

A natural question is whether the section subject to resonance ages faster than
the non-resonant one. The answer must separate the \emph{relative} effect of the
resonance from the \emph{absolute} stress level, and must keep the acoustic
vibration distinct from the mechanisms that actually govern the cable service
life.

\paragraph{Relative effect.}
Fatigue damage is driven by the amplitude of the cyclic stress. Resonance amplifies
the vibration amplitude, and therefore the bending stress at the cleats, by the
mechanical quality factor $Q$ (typically $10$--$50$). Since the fatigue life follows
a Basquin law
\begin{equation}
	N_f = N_{\mathrm{ref}}\left(\frac{\sigma_{\mathrm{ref}}}{\sigma_a}\right)^{k},
	\qquad k\approx 5\text{--}10,
	\label{eq:basquin}
\end{equation}
a factor $Q$ on the stress becomes a factor $Q^{k}$ on the damage rate. On this
basis the resonant section is, in relative terms, far more stressed than the
non-resonant one. This is the origin of the concern, and the reason the absolute
level must also be evaluated.

\paragraph{Absolute stress level.}
The decisive quantity is the actual vibration amplitude. The audible emission
estimated in Section~\ref{sec:onsite_summary} corresponds to a surface velocity of
about $\hat v \approx \SI{27}{\micro\meter\per\second}$ at resonance, i.e.\ a
transverse displacement amplitude
\begin{equation}
	w_0 = \frac{\hat v}{\omega_{sw}} \approx \SI{1}{\nano\meter}.
\end{equation}
The corresponding surface bending strain, with conductor radius $a$ and bending
wavenumber $k_B=2\pi/\lambda_B$, is
\begin{equation}
	\varepsilon_a \approx a\,k_B^{2}\,w_0 \sim 10^{-8},
	\label{eq:strain}
\end{equation}
giving a cyclic stress $\sigma_a = E\varepsilon_a$ of the order of a few
kilopascal, even including a cleat stress-concentration factor of $3$--$5$. The
copper fatigue strength in the very-high-cycle regime is of the order of tens of
megapascal; the applied stress is therefore three to four orders of magnitude
below the fatigue strength. Table~\ref{tab:fatigue} collects the estimate.

\begin{table}[h]
\centering
\caption{Order-of-magnitude estimate of the vibration-induced fatigue (rated load,
constant PWM, resonant span).}
\label{tab:fatigue}
\begin{tabular}{@{}lcl@{}}
\toprule
Quantity & Symbol & Value \\
\midrule
Displacement amplitude & $w_0$ & $\sim\SI{1}{\nano\meter}$ \\
Cyclic bending strain & $\varepsilon_a$ & $\sim 10^{-8}$ \\
Cyclic stress (incl.\ SCF) & $\sigma_a$ & $\sim\SI{5}{\kilo\pascal}$ \\
Copper fatigue strength (VHCF) & $\sigma_f$ & $\sim\SI{50}{\mega\pascal}$ \\
Stress margin & $\sigma_f/\sigma_a$ & $\sim 10^{4}$ \\
Cycles per year (50\% load) & $\dot n$ & $\sim\num{6e10}\ \mathrm{yr^{-1}}$ \\
Cycles to failure~\eqref{eq:basquin} ($k=7$) & $N_f$ & $\sim 10^{37}$ \\
\bottomrule
\end{tabular}
\end{table}

\paragraph{Cycle count and accumulated damage.}
At $f_{sw}=\SI{3.8}{\kilo\hertz}$ with a $50\%$ load duty, the excitation delivers
$\dot n\approx\num{6e10}$ cycles per year, i.e.\ about $\num{3e12}$ cycles over the
\num{50}-year design life --- a very-high-cycle regime. The Palmgren--Miner damage
accumulated over the design life $T_d$ is
\begin{equation}
	D = \frac{\dot n\,T_d}{N_f} \sim \frac{\num{3e12}}{10^{37}} \sim 10^{-25},
	\label{eq:miner}
\end{equation}
which is negligible by an enormous margin. The equivalent fatigue life imposed by
the acoustic vibration alone exceeds $10^{15}$ years and is not a service-life
constraint.

\paragraph{Aging threshold.}
To make the criterion explicit, define the fatigue life that the vibration
\emph{alone} would impose, $L_{\mathrm{fat}}=N_f/\dot n$, and the resulting aging
fraction over the design life,
\begin{equation}
	\varphi = \frac{T_d}{L_{\mathrm{fat}}} = D.
\end{equation}
The following bands are adopted:
\begin{equation}
	\boxed{\;
	\begin{aligned}
	&\varphi < 0.02 \;\;(\text{1 yr in 50}) && \Rightarrow\;\; \text{negligible},\\
	&0.02 \le \varphi \le 0.20 && \Rightarrow\;\; \text{monitor},\\
	&\varphi > 0.20 \;\;(\text{10 yr in 50}) && \Rightarrow\;\; \text{significant}.
	\end{aligned}
	\;}
\end{equation}
The present estimate, $\varphi\sim10^{-25}$, falls in the negligible band with a
margin of many orders of magnitude.

\paragraph{Qualifications.}
Three points must accompany this result in order not to misread it:
\begin{itemize}[leftmargin=*]
	\item[--] \textbf{Acoustic emission is not a damage indicator.} The audible tone
	is a comfort/acceptability issue, not a measure of fatigue. In fact the
	$2f_1$ ($\SI{100}{\hertz}$) force is about seven times larger than the
	switching-frequency force yet is inaudible: if a fatigue driver ever existed,
	it would be the largest force acting on the most compliant mode, independent of
	its audibility. The noise must not be used as a proxy for damage.
	\item[--] \textbf{The service life is governed by other mechanisms.} The
	\num{50}-year life is set by thermal aging of the insulation
	(Arrhenius law, IEC~60216), by electrical aging (partial discharge, water
	treeing in the XLPE), and by thermo-mechanical fatigue at the cleats driven by
	the \emph{load cycle} (daily/seasonal expansion, a few cycles per day --- not
	the kilohertz vibration). The PWM-induced resonance is negligible relative to
	these.
	\item[--] \textbf{The only vibration-related watch item} is the possible
	progressive loosening or fretting of cleats, terminations and clamps. At a
	displacement amplitude of order \SI{1}{\nano\meter} the motion is well below the
	micrometre slip threshold of fretting, so this is a periodic-inspection item,
	not a life limit.
\end{itemize}

\paragraph{Confirmation by measurement.}
Because the fatigue damage scales as $Q^{k}$, the absolute figure is sensitive to
the true strain at the cleat, which has here been inferred acoustically. A direct
measurement closes the question:
\begin{enumerate}[leftmargin=*]
	\item Instrument a \SI{120}{\meter} span at the point of maximum curvature (at a
	cleat) with a strain gauge, or with an accelerometer converted to strain through
	the local mode shape.
	\item Record under rated load with constant PWM (the worst, tonal case), and
	extract the cyclic strain amplitude $\varepsilon_a$ at $f_{sw}\pm m f_1$.
	\item Repeat off resonance to bound $Q$.
	\item Insert the measured $\varepsilon_a$ into \eqref{eq:basquin}--\eqref{eq:miner}
	with the cable-material S--N data and compare $\varphi$ against the threshold.
\end{enumerate}
This converts the assessment from an estimate into a verified result.
